% GENERATED FILE. Edit canonical lessons, then run npm run generate:books.
% canonical-source-hash: fnv1a64:2c0c9051d8dca65a
% canonical-lessons: RU-C505-rybalka, RU-C505-udochka, RU-C505-kryuchok, RU-C505-klevat, RU-C505-ulov

\chapter{Gone Fishing}
\label{ch:ru-gone-fishing}
\hlchaptermodality{505}

\begin{chapteropening}
\textbf{By the end of this chapter:} \emph{I can talk about a fishing trip.}

Complete the last lesson of chapter 505: I can talk about a fishing trip.
\end{chapteropening}

\section[\texorpdfstring{\ru{рыбалка} (rybálka)}{рыбалка (rybálka)}]{\ru{рыбалка} (rybálka) — fishing}

\label{lesson:RU-C505-rybalka}



\begin{quote}
Before the new one: say the Russian for a pack of cards, then the Russian for to cheat.
\end{quote}

\subsection*{You'll want to know: \ru{рыбалка}}
\textbf{\ru{рыбалка}} (\emph{rybálka}) — \textquotedblleft{}fishing\textquotedblright{}.

\begin{quote}
\ru{В} \ru{субботу} \ru{мы} \ru{с} \ru{папой} \ru{едем} \ru{на} \ru{рыбалку}.
\end{quote}

\begin{quote}
\emph{V subbótu my s pápoy yédem na rybálku.}
\end{quote}

\begin{quote}
On Saturday Dad and I are going fishing.
\end{quote}

Feminine. \emph{\ru{Ехать} \ru{на} \ru{рыбалку}} — to go fishing.

\begin{cousinweb}
\emph{\ru{рыбалка}} is built on \emph{\ru{рыба}} 'fish', through \emph{\ru{рыбак}} 'fisherman'; EtymDB has no entry.
\end{cousinweb}

\begin{grammarlens}[title={What you've built}]
One more word for telling, explaining and describing at B1.
\end{grammarlens}

\subsection*{Your turn}
\begin{itemize}
\raggedright
  \item \emph{Say it:} \emph{rybálka}
  \item \emph{Say it:} \emph{rybálka}, once more
  \item \emph{Say it:} a sentence of your own with \emph{rybálka}
  \item \emph{Recall:} say the Russian for a pack of cards, then the Russian for to cheat, then say \emph{rybálka} again
\end{itemize}

\begin{tcolorbox}[breakable,colback=teal!4,colframe=teal!35!black,title={Before you move on}]
What does \ru{рыбалка} mean? (\textquotedblleft{}Fishing\textquotedblright{}.) Say it once more.
\end{tcolorbox}
\section[\texorpdfstring{\ru{удочка} (údochka)}{удочка (údochka)}]{\ru{удочка} (údochka) — a fishing rod}

\label{lesson:RU-C505-udochka}



\begin{quote}
Before the new one: say the Russian for to cheat, then the Russian for fishing.
\end{quote}

\subsection*{You'll want to know: \ru{удочка}}
\textbf{\ru{удочка}} (\emph{údochka}) — \textquotedblleft{}a fishing rod\textquotedblright{}.

\begin{quote}
\ru{Он} \ru{взял} \ru{удочку} \ru{и} \ru{пошёл} \ru{к} \ru{озеру}.
\end{quote}

\begin{quote}
\emph{On vzyal údochku i poshyól k ózeru.}
\end{quote}

\begin{quote}
He took his fishing rod and went to the lake.
\end{quote}

Feminine. A diminutive of the old word \emph{\ru{уда}}.

\begin{cousinweb}
\emph{\ru{удочка}} is a diminutive of the old word \emph{\ru{уда}} 'a fishing line'; EtymDB has no entry, and \emph{\ru{уда}} is usually traced to Proto-Slavic \emph{ǫda}.
\end{cousinweb}

\begin{grammarlens}[title={What you've built}]
One more word for telling, explaining and describing at B1.
\end{grammarlens}

\subsection*{Your turn}
\begin{itemize}
\raggedright
  \item \emph{Say it:} \emph{údochka}
  \item \emph{Say it:} \emph{údochka}, once more
  \item \emph{Say it:} a sentence of your own with \emph{údochka}
  \item \emph{Recall:} say the Russian for to cheat, then the Russian for fishing, then say \emph{údochka} again
\end{itemize}

\begin{tcolorbox}[breakable,colback=teal!4,colframe=teal!35!black,title={Before you move on}]
What does \ru{удочка} mean? (\textquotedblleft{}A fishing rod\textquotedblright{}.) Say it once more.
\end{tcolorbox}
\section[\texorpdfstring{\ru{крючок} (kryuchók)}{крючок (kryuchók)}]{\ru{крючок} (kryuchók) — a hook}

\label{lesson:RU-C505-kryuchok}



\begin{quote}
Before the new one: say the Russian for fishing, then the Russian for a fishing rod.
\end{quote}

\subsection*{You'll want to know: \ru{крючок}}
\textbf{\ru{крючок}} (\emph{kryuchók}) — \textquotedblleft{}a hook\textquotedblright{}.

\begin{quote}
\ru{На} \ru{крючок} \ru{попалась} \ru{маленькая} \ru{рыбка}.
\end{quote}

\begin{quote}
\emph{Na kryuchók popalás málenkaya rýbka.}
\end{quote}

\begin{quote}
A little fish got caught on the hook.
\end{quote}

Masculine; genitive \ru{крючка}. Also a coat hook.

\begin{cousinweb}
\emph{\ru{крючок}} is a diminutive of \emph{\ru{крюк}} 'a hook'; EtymDB has no entry.
\end{cousinweb}

\begin{grammarlens}[title={What you've built}]
One more word for telling, explaining and describing at B1.
\end{grammarlens}

\subsection*{Your turn}
\begin{itemize}
\raggedright
  \item \emph{Say it:} \emph{kryuchók}
  \item \emph{Say it:} \emph{kryuchók}, once more
  \item \emph{Say it:} a sentence of your own with \emph{kryuchók}
  \item \emph{Recall:} say the Russian for fishing, then the Russian for a fishing rod, then say \emph{kryuchók} again
\end{itemize}

\begin{tcolorbox}[breakable,colback=teal!4,colframe=teal!35!black,title={Before you move on}]
What does \ru{крючок} mean? (\textquotedblleft{}A hook\textquotedblright{}.) Say it once more.
\end{tcolorbox}
\section[\texorpdfstring{\ru{клевать} (klevát)}{клевать (klevát)}]{\ru{клевать} (klevát) — to bite (of fish); to peck}

\label{lesson:RU-C505-klevat}



\begin{quote}
Before the new one: say the Russian for a fishing rod, then the Russian for a hook.
\end{quote}

\subsection*{You'll want to know: \ru{клевать}}
\textbf{\ru{клевать}} (\emph{klevát}) — \textquotedblleft{}to bite (of fish); to peck\textquotedblright{}.

\begin{quote}
\ru{Сегодня} \ru{рыба} \ru{плохо} \ru{клюёт}.
\end{quote}

\begin{quote}
\emph{Segódnya rýba plókho klyuyót.}
\end{quote}

\begin{quote}
The fish aren't biting well today.
\end{quote}

Imperfective; present \ru{клюю}, \ru{клюёшь} (-\ru{ева}- becomes -\ru{ю}-). Perfective \ru{клюнуть}.

\begin{cousinweb}
\emph{\ru{клевать}} is inherited from Proto-Slavic \emph{kľ\ru{ь}vati} 'to peck, to bite' (EtymDB).
\end{cousinweb}

\begin{grammarlens}[title={What you've built}]
One more word for telling, explaining and describing at B1.
\end{grammarlens}

\subsection*{Your turn}
\begin{itemize}
\raggedright
  \item \emph{Say it:} \emph{klevát}
  \item \emph{Say it:} \emph{klevát}, once more
  \item \emph{Say it:} a sentence of your own with \emph{klevát}
  \item \emph{Recall:} say the Russian for a fishing rod, then the Russian for a hook, then say \emph{klevát} again
\end{itemize}

\begin{tcolorbox}[breakable,colback=teal!4,colframe=teal!35!black,title={Before you move on}]
What does \ru{клевать} mean? (\textquotedblleft{}To bite (of fish); to peck\textquotedblright{}.) Say it once more.
\end{tcolorbox}
\section[\texorpdfstring{\ru{улов} (ulóv)}{улов (ulóv)}]{\ru{улов} (ulóv) — a catch (of fish)}

\label{lesson:RU-C505-ulov}



\begin{quote}
Before the new one: say the Russian for a hook, then the Russian for to bite.
\end{quote}

\subsection*{You'll want to know: \ru{улов}}
\textbf{\ru{улов}} (\emph{ulóv}) — \textquotedblleft{}a catch (of fish)\textquotedblright{}.

\begin{quote}
\ru{Рыбаки} \ru{были} \ru{довольны} \ru{уловом}.
\end{quote}

\begin{quote}
\emph{Rybáki býli dovólny ulóvom.}
\end{quote}

\begin{quote}
The fishermen were pleased with their catch.
\end{quote}

Masculine.

\begin{cousinweb}
\emph{\ru{улов}} is built on \emph{\ru{у}-} and the root of \emph{\ru{ловить}} 'to catch'; EtymDB has no entry.
\end{cousinweb}

\begin{grammarlens}[title={What you've built}]
One more word for telling, explaining and describing at B1.
\end{grammarlens}

\subsection*{Your turn}
\begin{itemize}
\raggedright
  \item \emph{Say it:} \emph{ulóv}
  \item \emph{Say it:} \emph{ulóv}, once more
  \item \emph{Say it:} a sentence of your own with \emph{ulóv}
  \item \emph{Recall:} say the Russian for a hook, then the Russian for to bite, then say \emph{ulóv} again
\end{itemize}

\begin{tcolorbox}[breakable,colback=teal!4,colframe=teal!35!black,title={Before you move on}]
What does \ru{улов} mean? (\textquotedblleft{}A catch (of fish)\textquotedblright{}.) Say it once more.
\end{tcolorbox}
